As a global optimization scheme, the \ac{de} is implemented, using \textit{scipy}. Its strength lies in generating multiple experimental vectors (mutations) to search in multiple directions, then retaining the parameters (crossover) that yield the lowest objective function. This allows the algorithm in the early state to explore multiple areas. The population size of the \ac{de} determines the number of possible directions, while a high value increases the exploration and simultaneously decreases the convergence rate.

\subsubsection{Differential Evolution (DE) Algorithm}
\cite{journals/jgo/StornP97}
Differential Evolution is a parallel direct search method designed for global optimization over continuous spaces. In the context of optimizing the proposed model, DE iteratively updates a population of $NP$ $D$-dimensional parameter vectors $\mathbf{x}_{i,G}$ ($i=1, 2, \dots, NP$) through three fundamental operators:

\begin{itemize}
    \item \textbf{Mutation:} For each target vector, a mutant vector $\mathbf{v}_{i,G+1}$ is generated by adding the weighted difference between two random population vectors to a third vector:
    \begin{equation}
        \mathbf{v}_{i,G+1} = \mathbf{x}_{r_1,G} + F \cdot (\mathbf{x}_{r_2,G} - \mathbf{x}_{r_3,G})
    \end{equation}
    where $r_1, r_2, r_3$ are mutually different random indices, and $F \in [0, 2]$ is the scaling factor that controls the amplification of the differential variation.
    
    \item \textbf{Crossover:} To increase diversity, a trial vector $\mathbf{u}_{i,G+1}$ is formed by mixing parameters from the target vector and the mutant vector based on a crossover constant $CR \in [0, 1]$:
    \begin{equation}
        u_{ji,G+1} = 
        \begin{cases} 
        v_{ji,G+1} & \text{if } (randb(j) \le CR) \text{ or } j = rnbr(i) \\ 
        x_{ji,G} & \text{otherwise} 
        \end{cases}
    \end{equation}
    where $rnbr(i)$ is a randomly chosen index to ensure the trial vector receives at least one parameter from the mutant vector.
    
    \item \textbf{Selection:} A greedy criterion is applied where the trial vector replaces the target vector in the next generation only if it yields a lower cost function value:
    \begin{equation}
        \mathbf{x}_{i,G+1} = 
        \begin{cases} 
        \mathbf{u}_{i,G+1} & \text{if } f(\mathbf{u}_{i,G+1}) \le f(\mathbf{x}_{i,G}) \\ 
        \mathbf{x}_{i,G} & \text{otherwise} 
        \end{cases}
    \end{equation}
\end{itemize}



\subsubsection{Code Implementation}
Due to its minimal control variables, ease of use, and self-organizing nature, DE effectively optimizes complex models that involve non-differentiable or multi-modal landscapes.In summary, DE is specifically designed to solve nonlinear, even non-differentiable continuous-space optimization problems, making it highly suitable for complex environmental models like HBV model. Through the objective function we set, it iteratively finds the input parameters that minimize this function value. In Assignment 1, Our optimization goal is to find a set of HBV parameters that minimizes the "error" between the runoff simulated by the HBV model and the observed actual runoff. The "error function" (NSE or RMSE) serves as the objective function optimized by the DE algorithm.
In the code, we vary the population size parameter to obtain results corresponding to different population sizes, enabling us to utilize parameter variations for comparative analysis. Simultaneously, we analyze parameter sensitivity by altering the boundaries of parameters within the HBV model.


The main code is presented in table \ref{tab:code_01}. The input required to run the HBV001a model is collected as objects in the \textit{recorder} class. This in turn serves as input (\textit{func}, mandatory) to the \ac{de} scheme and returns the \ac{ofv} for further evaluation. Further, the \textit{recorder} class records information about the \ac{mp} 's development during the optimization process for later evaluation. As efficiency metrics, the objective function of the \ac{nse} is applied to achieve a good fit for the high flows. Only two inputs are mandatory for \ac{de}: the already managed \textit{func} and the \ac{mp} boundaries. The number of generations (\textit{maxiter}) and the population size (\textit{popsize}) define the maximum number of model runs used for the optimization process: \textit{maxiter} * \textit{popsize} * number of \ac{mp}. \textit{atol} is used to generate a stop criterion based on the number of significant decimals. For the \ac{nse} value, results ranging from $0.75 < \mathrm{NSE} \le 1$ are considered very good, and $0.65 < \mathrm{NSE} \le 1$0.75 as good \autocite{Moriasi.2007}. In conclusion, three significant decimals are sufficient. 

\begin{table}[H]
\centering
\begin{tabular}{|l|}
\hline
create main   directory path \\ \hline
read input: time series (*.csv) \\ \hline
define \ac{mp} boundary \\ \hline
check correct \ac{mp} order for HBV001a \\ \hline
define generation and population size for   optimizing algorithms  \\ \hline
\begin{tabular}[c]{@{}l@{}}creat   "recorder" class: \\ 
- input: temperature, precipitation, evapotranspiration, \\      
Conversion constant from \si{mm/h} to \si{m^3/s}, \\
time steps, observed discharge, model, and population size\\      
- call \underline{function}: runs HBV001a model \\
- output: \ac{ofv} of the \ac{nse}
\end{tabular} \\ \hline
\begin{tabular}[c]{@{}l@{}}
\underline{function}: \textbf{Differential Evolution}\\      
- input: recoder, \ac{mp} boundaries, generation size, population size\\      
- determination: strategy = 'best1bin', atol=1e-3, polish=False\\      
- output: best \ac{mp} set, and related \ac{ofv}
\end{tabular} \\ \hline
create output directory \\ \hline
generate plots \\ \hline
\end{tabular}
\caption{Optimization: main code structure}
\label{tab:code_01}
%\end{table}
\end{table}

To turn off some of the processes, a separate script utilizes the provided aa\_run\_hbv001.py to run the model with the altered \ac{mp} sets, plot the hydrograph and other process values, and record the \ac{nse} values. To turn off the upper reservoir (urr), the fast responding outflows from the urr model are reduced to zero (\ac{mp}: urr\_tdr and urr\_ndr), and the ratio of the percolation transfers is set to 1 (\ac{mp}: urr\_ulc). This way, all water is transferred from the upper reservoir to the lower reservoir (lrr). For the lower reservoir, the inflow (\ac{mp}: urr\_ulc) and the outflow (\ac{mp}: urr\_dre) are set to zero. Although it would have been sufficient to just turn off the inflow since it is the last model in the cascade. The last two processes turned off are the snow module, with the temperature threshold for snow set to \SI{-100}{\degreeCelsius}, and the \ac{gw} module. Concerning the \ac{gw} module, the flow that is entering the river in the catchment (\ac{mp}: Irr\_uct) and the flow leaving the catchment unnoticed (\ac{mp}: Lrr\_lct) are set to zero.  

\begin{comment}   
Was ist gegeben 
Welche Konzepte Benutzen wir NSE, logNSE, DE - Konzepte dahinter 
Wie implemtniert (Diagramm?) 
Besonderheit - erkenntnisse aus späteren Übungen

• Implement various efficiency metrics
− NSE, LnNSE 
• Implement the objective function (OF)
− Accepts model parameters
− Accepts model and related data
− Runs model
− Computes efficiency metrics
− Returns objective function value (OFV)
• Implement the optimization scheme
\end{comment}